\cotefile{157-2}
\scanpage{147}
\mspage{92}
ce qui \unsure{résultera} aussitôt du corollaire :

\uline{Corollaire} \unsure{Soient|Soit} un diagramme dans $\mathcal M_{/S}$
\diagram{\begin{tikzcd}[ampersand replacement=\&]
X\arrow[r,"i"] \& C\arrow[r,shift left=.6ex,"u_1"]\arrow[r,shift right=.6ex,"u_2"'] \& C'\text{,}
\end{tikzcd}}
avec $C$ et $C'$ dans $\underline{\mathrm{Fib}}\,S$ (i.e. $C\longrightarrow S$ et $C'\longrightarrow S$ dans $\mathrm{Fib}$), et $i\in\mathrm{Cof}_0$\inserted{ et $\unsure{i|i'}\in W$}. Si $u_1i=u_2i$, alors $u_1$ et $u_2$ définissent la $\eqsim$ flèche de $(\underline{\mathrm{Fib}}\,S)W_S^{-1}$.

\uline{Dém} Soit
\[
u=(u_1,u_2):C\longrightarrow C'\times C'\text{,}
\]
qui rend commutatif par hyp. le diagramme
\diagram{\begin{tikzcd}[ampersand replacement=\&]
X\arrow[r,"w"]\arrow[d,"i"'] \& C'\arrow[d,"\mathrm{diag}_{C'}"]\\
C\arrow[r,"u"'] \& C'\times_S C'\text{,}
\end{tikzcd}}
où $w=u_1i=u_2i$. Factorisons $\mathrm{diag}_{C'}$ en
\diagram{\begin{tikzcd}[ampersand replacement=\&]
C'\arrow[r,"i'"] \& \overline{C'}\arrow[r,"\delta'"] \& C'\times_S C'\text{,}\quad\text{avec }\delta'\in\mathrm{Fib}\text{,}
\end{tikzcd}}
\deleted{\unsure{pas|par}} et $i'\in W$. On trouve un carré comm.
\diagram{\begin{tikzcd}[ampersand replacement=\&]
X\arrow[r,"w'=i'w"]\arrow[d,"i"'] \& \overline{C'}\arrow[d,"\delta'"] \& {}\\
C\arrow[ur,dashed,"\overline u"]\arrow[r,"u"] \& C'\times_S C'\arrow[r,shift left=.6ex,"\mathrm{pr}_1"]\arrow[r,shift right=.6ex,"\mathrm{pr}_2"'] \& C'
\end{tikzcd}}
avec $i\in\mathrm{Cof}_0$, $\delta'\in\mathrm{Fib}$, d'où l'existence de la flèche en pointillés, rendant les deux triangles commutatifs. On veut prouver $\mathrm{pr}_1\circ u=\mathrm{pr}_2\circ u$ dans $(\underline{\mathrm{Fib}}\,S)W_S^{-1}$, comme $u=\delta'\overline u$ il suffit de prouver que l'on a $\mathrm{pr}_1\delta'=\unsure{p_2|\mathrm{pr}_2}\delta'$ dans $(\underline{\mathrm{Fib}}\unsure{,}S)W_S^{-1}$.
% archive sheet number: 146
